%% ============================================================
%% MODELO DE ATIVIDADE · ENSAIO EXPERIMENTAL E VALIDAÇÃO
%%
%% Copie este arquivo para atividades/ (atividade-02.tex, por exemplo),
%% inclua-o no relatorio-aacc.tex com \input{atividades/atividade-02}
%% e escreva. Cada \preencher{…} é uma pendência: troque-o pelo seu
%% texto. Cada \figuraprovisoria também: troque-a pela figura de verdade.
%%
%% Uma atividade é UMA entrega sua, com começo, meio e fim. "Participei
%% do projeto X por dois anos" não é uma atividade: são várias. Divida.
%% ============================================================
\begin{atividade}{Ensaio de V na bancada: montagem, medições e critério de aceite}{
  tipo         = ensaio,
  modalidade   = {},   % a chave da modalidade pedida: EXT, PES, COMP, VPC, EVT, PRD, ORG, ENS, GES, CUR, OUT
  alternativa  = {},   % a modalidade alternativa, se o Colegiado não aceitar a primeira
  horas        = {},   % as horas desta atividade, com base nos registros (Parte IV)
  inicio       = {},   % AAAA-MM
  fim          = {},   % AAAA-MM, ou: em andamento
  projeto      = {},
  papel        = {},   % Responsável, Corresponsável ou Colaborador(a)
  supervisor   = {},   % quem confirma: nome e cargo (o supervisor do projeto, o gerente)
  registros    = {},   % os códigos no SOMA, com ponto e vírgula: NBL-14; NRO-PRO-003-7
  competencias = {},   % as competências das DCN que ela desenvolveu (I a VIII): I, III, V
  disciplinas  = {},   % as disciplinas do seu curso ligadas a ela: ELE075 Eletrônica; ELE067 Circuitos
}

\begin{contexto}
\preencher{O que foi ensaiado, contra que requisito ou hipótese, e por que o ensaio era necessário.}
\end{contexto}

\begin{contribuicao}
\preencher{O que você planejou, montou, mediu e analisou.}
\end{contribuicao}

\begin{desenvolvimento}
\fichatecnica{
  protocolo    = {},   % Protocolo ou plano de ensaio
  instrumentos = {},   % Instrumentos (e calibração)
  local        = {},   % Local
  amostra      = {},   % Amostras ou participantes — opcional
  etica        = {},   % Aprovação ética (CAAE), se houve pessoas — opcional
  relatorio    = {},   % Relatório de execução (NRO-PRO-003)
}

\preencher{A montagem experimental, os instrumentos (com a resolução, a incerteza e a calibração), o procedimento, o que foi controlado e o que foi repetido. Com pessoas: o CAAE, o consentimento e a supervisão profissional.}

\begin{figure}[htbp]
  \centering
  \figuraprovisoria{A montagem experimental (foto ou diagrama), com os instrumentos identificados.}
  % \includegraphics[width=0.9\linewidth]{figuras/montagem.pdf}
  \caption{Montagem do ensaio de V na bancada do LABBIO.}
  \fonte{Elaborado pelo autor.}
  \label{fig:montagem}
\end{figure}
\end{desenvolvimento}

\begin{resultados}
\preencher{Tabelas e gráficos com unidades e incertezas; a análise; o critério de aceite, atendido ou não. O relatório de execução (NRO-PRO-003).}
\end{resultados}

\begin{evidencias}
% Uma linha por evidência: \evidencia{tipo}{identificador}{o que mostra}{onde conferir}
% \evidencia{Arquivo controlado}{NRO-PRO-003-<PN>}{O relatório de execução do ensaio}{SOMA › Arquivos}
% \evidencia{Dados}{<projeto>/ensaios/<data>.csv}{As medições brutas}{Drive da equipe (acesso a pedido)}
\end{evidencias}

\begin{aprendizados}
\preencher{O que a medição de um sistema real ensinou (e as disciplinas de laboratório, de instrumentação e de estatística que ela exigiu).}
\end{aprendizados}
\end{atividade}
